\documentclass{article}
\usepackage{rkdiagram}

\begin{document}
	
	\section*{Step 2: main-clause baseline layout}
	
	\subsection*{Subject + Verb}
	
	\begin{rksentence}
		\rksubject{dog}
		\rkverb{ran}
	\end{rksentence}
	
	\subsection*{Subject + Verb + Object}
	
	\begin{rksentence}
		\rksubject{dog}
		\rkverb{chased}
		\rkobject{cat}
	\end{rksentence}
	
	\subsection*{Predicate adjective (backward-slanting divider)}
	
	\begin{rksentence}
		\rksubject{dog}
		\rkverb{is}
		\rkpredadj{happy}
	\end{rksentence}
	
	\subsection*{Longer words — auto-width scaling}
	
	\begin{rksentence}
		\rksubject{elephant}
		\rkverb{remembered}
		\rkobject{everything}
	\end{rksentence}
	
	% --------------------------------------------------------
	\section*{Step 3: modifier commands}
	
	\subsection*{Single adjective on subject}
	% "The dog ran" — article hangs below subject
	
	\begin{rksentence}
		\rksubject{dog}
		\rkverb{ran}
		\rkadj{The}
	\end{rksentence}
	
	\subsection*{Single adverb on verb}
	% "dog ran quickly"
	
	\begin{rksentence}
		\rksubject{dog}
		\rkverb{ran}
		\rkadv{quickly}
	\end{rksentence}
	
	\subsection*{Stacked modifiers — two adjectives on subject}
	% "The big dog ran"
	
	\begin{rksentence}
		\rksubject{dog}
		\rkverb{ran}
		\rkadj{The}
		\rkadj{big}
	\end{rksentence}

	\subsection*{Source-order stacking on one slot}
	% modifiers should appear on separate slanted lines in code order

	\begin{rksentence}
		\rksubject{dog}
		\rkverb{ran}
		\rkadj{The}
		\rkadj{big}
		\rkadj{brown}
	\end{rksentence}
	
	\subsection*{Modifiers on multiple slots}
	% "The big dog ran very quickly"
	
	\begin{rksentence}
		\rksubject{dog}
		\rkverb{ran}
		\rkadj{The}
		\rkadj{big}
		\rkadv{very}
		\rkadv{quickly}
	\end{rksentence}
	
	\subsection*{Adjective on object}
	% "dog chased the small cat" — two adjectives on object
	
	\begin{rksentence}
		\rksubject{dog}
		\rkverb{chased}
		\rkobject{cat}
		\rkadj[attach=object]{the}
		\rkadj[attach=object]{small}
	\end{rksentence}
	
	\subsection*{Full sentence with predicate adjective and modifiers}
	% "The old dog is very happy"
	
	\begin{rksentence}
		\rksubject{dog}
		\rkverb{is}
		\rkpredadj{happy}
		\rkadj{The}
		\rkadj{old}
		\rkadv[attach=predadj]{very}
	\end{rksentence}
	
	\subsection*{Inline use}
	
	The sentence
	\begin{rksentence}
		\rksubject{dog}
		\rkverb{ran}
		\rkadj{The}
	\end{rksentence}
	is a classic example.
	
	% --------------------------------------------------------
	\section*{Step 4: prepositional phrases}
	
	\subsection*{Single prep phrase on verb}
	% "dog ran in park"
	
	\bigskip
	\begin{rksentence}
		\rksubject{dog}
		\rkverb{ran}
		\rkprep[attach=verb]{in}
		\rkprepobj{park}
	\end{rksentence}
	
	\subsection*{Prep phrase on subject}
	% "dog in yard ran"
	
	\bigskip
	\begin{rksentence}
		\rksubject{dog}
		\rkverb{ran}
		\rkprep[attach=subject]{in}
		\rkprepobj{yard}
	\end{rksentence}
	
	\subsection*{Two prep phrases on the same slot}
	% "dog ran in park on trail" — both preps on verb, stacking below each other
	
	\bigskip
	\begin{rksentence}
		\rksubject{dog}
		\rkverb{ran}
		\rkprep[attach=verb]{in}
		\rkprepobj{park}
		\rkprep[attach=verb]{on}
		\rkprepobj{trail}
	\end{rksentence}
	
	\subsection*{Modifiers and a prep phrase together}
	% "The big dog ran quickly in the park"
	
	\bigskip
	\begin{rksentence}
		\rksubject{dog}
		\rkverb{ran}
		\rkadj{The}
		\rkadj{big}
		\rkadv{quickly}
		\rkprep[attach=verb]{in}
		\rkprepobj{park}
	\end{rksentence}
	
	\subsection*{Prep phrases on both subject and verb}
	% "The dog in the yard ran in the park"
	
	\bigskip
	\begin{rksentence}
		\rksubject{dog}
		\rkverb{ran}
		\rkadj{The}
		\rkprep[attach=subject]{in}
		\rkprepobj{yard}
		\rkprep[attach=verb]{in}
		\rkprepobj{park}
	\end{rksentence}
	
	% --------------------------------------------------------
	\section*{Step 5: compound structures}
	
	\subsection*{Compound subject}
	% "Dogs and cats ran"
	
	\bigskip
	\begin{rksentence}
		\rkcompound[coord=and]{dogs}{cats}
		\rkverb{ran}
	\end{rksentence}
	
	\subsection*{Compound subject with object}
	% "Dogs and cats chased the bird"
	
	\bigskip
	\begin{rksentence}
		\rkcompound[coord=and]{dogs}{cats}
		\rkverb{chased}
		\rkobject{bird}
		\rkadj[attach=object]{the}
	\end{rksentence}
	
	\subsection*{Compound predicate}
	% "dog ran and barked"
	
	\bigskip
	\begin{rksentence}
		\rksubject{dog}
		\rkcompound[role=verb, coord=and]{ran}{barked}
	\end{rksentence}
	
	\subsection*{Compound predicate with object}
	% "dog chased and caught the cat"
	
	\bigskip
	\begin{rksentence}
		\rksubject{dog}
		\rkcompound[role=verb, coord=and]{chased}{caught}
		\rkobject{cat}
		\rkadj[attach=object]{the}
	\end{rksentence}
	
	\subsection*{Compound subject with or}
	% "dogs or cats ran"
	
	\bigskip
	\begin{rksentence}
		\rkcompound[coord=or]{dogs}{cats}
		\rkverb{ran}
	\end{rksentence}
	
	\subsection*{Compound subject with dependent verb phrase}
	% "dogs and cats ran quickly in the yard"
	
	\bigskip
	\begin{rksentence}
		\rkcompound[coord=and]{dogs}{cats}
		\rkverb{ran}
		\rkadv{quickly}
		\rkprep[attach=verb]{in}
		\rkprepobj{yard}
	\end{rksentence}
	
	% --------------------------------------------------------
	\section*{Step 6: experimental structures}
	
	This section keeps examples for partially implemented commands.
	The output here is useful for development, but it is not yet a
	fidelity check against full traditional Reed--Kellogg notation.
	
	\subsection*{Appositive on subject}
	% current implementation places the appositive after the noun slot
	
	\bigskip
	\begin{rksentence}
		\rksubject{dog}
		\rkappositive[attach=subject]{terrier}
		\rkverb{ran}
	\end{rksentence}
	
	\subsection*{Appositive on object}
	
	\bigskip
	\begin{rksentence}
		\rksubject{dog}
		\rkverb{chased}
		\rkobject{cat}
		\rkappositive[attach=object]{tabby}
	\end{rksentence}
	
	\subsection*{Gerund as subject}
	
	\bigskip
	\begin{rksentence}
		\rkgerund[role=subject]{Running}
		\rkverb{is}
		\rkpredadj{fun}
	\end{rksentence}
	
	\subsection*{Gerund as object}
	
	\bigskip
	\begin{rksentence}
		\rksubject{dog}
		\rkverb{loves}
		\rkgerund[role=object]{running}
	\end{rksentence}
	
\end{document}
